\documentclass[12pt,a4paper]{article}
\usepackage{amsmath,amssymb,amsthm}
\newtheorem{theorem}{Theorem}
\newtheorem{lemma}[theorem]{Lemma}
\theoremstyle{definition}
\newtheorem{definition}{Definition}
\title{A Few Remarks on Geometric Series}
\author{Tobias Reinholt}
\date{June 14, 2023}
\begin{document}
\maketitle

\section{Definitions}
\begin{definition}\label{def:geo}
A \emph{geometric series} is a series of the form $\sum_{k=0}^{\infty} r^k$,
where $r$ is a real number.
\end{definition}

The partial sums of such a series can be computed in closed form. For
$r \neq 1$ we have
\begin{equation}\label{eq:partial}
S_n = \sum_{k=0}^{n} r^k = \frac{1 - r^{n+1}}{1 - r}.
\end{equation}

\begin{lemma}\label{lem:partial}
Equation~\eqref{eq:partial} holds for every integer $n \geq 0$ and every
real $r \neq 1$.
\end{lemma}
\begin{proof}
For $n = 0$ both sides equal $1$. Suppose the claim holds for $n$. Then
\[
S_{n+1} = S_n + r^{n+1} = \frac{1 - r^{n+1}}{1 - r} + r^{n+1}
        = \frac{1 - r^{n+2}}{1 - r},
\]
which is the claim for $n+1$.
\end{proof}

\section{Convergence}
\begin{theorem}\label{thm:conv}
If $|r| < 1$ then $\sum_{k=0}^{\infty} r^k = \dfrac{1}{1-r}$.
\end{theorem}
\begin{proof}
By Lemma~\ref{lem:partial}, $S_n = (1 - r^{n+1})/(1-r)$. Since $|r|<1$, the
term $r^{n+1}$ tends to $0$ as $n \to \infty$.
\end{proof}

An immediate consequence, referred to on page~\pageref{thm:conv}, is that
\begin{align}
\sum_{k=1}^{\infty} 2^{-k} &= 1, \\
\sum_{k=0}^{\infty} \left(-\tfrac{1}{3}\right)^k &= \tfrac{3}{4}.
\end{align}

\section*{An Integral Analogue}
The continuous counterpart replaces the sum by an integral:
$\int_0^\infty e^{-\lambda t}\,dt = 1/\lambda$ for $\lambda > 0$. Combining
the two, one finds for $\lambda > 0$ that
\[
\int_0^{\infty} \sum_{k=0}^{\infty} e^{-(k+1)\lambda t}\,dt
  = \sum_{k=1}^{\infty} \frac{1}{k\lambda},
\]
a harmonic series, which diverges. Theorem~\ref{thm:conv} therefore cannot be
applied naively to the integrand at $t = 0$.
\end{document}
